\originalchapter{染色与 Tutte 多项式}\label{ch:polynomials}
\sourceof{18.315 L9--12; 计数、符号消去和活动性证明补全}
\originalsection{删除与收缩}
\begin{definition}
对非负整数 $q$, 令 $P_G(q)$ 为图 $G$ 的正常 $q$ 色顶点染色数. 缩边 $G/e$ 把非自环边 $e=uv$ 的两端识别为一点, 其余边随端点一起映射; 允许出现平行边和自环. 平行边不增加顶点染色约束, 自环使正常染色数为0.
\end{definition}
\begin{theorem}\label{thm:chromaticdel}
对非自环边 $e$, $P_G(q)=P_{G-e}(q)-P_{G/e}(q)$. 因而 $P_G$ 是多项式; 无自环 $n$ 顶点图的首项为 $q^n$.
\end{theorem}
\begin{proof}
删边图的染色分为端点异色和同色两类. 前者是原图染色; 后者与缩边图的染色一一对应. 无边图有 $q^n$ 种染色. 对边数归纳, 删除项顶点数不变, 收缩项少一个顶点, 得多项式与首项. 对平行边可先删除重复约束, 不影响结论.
\end{proof}
树的染色多项式为 $q(q-1)^{n-1}$: 根有 $q$ 种选择, 按父先子后, 每点只需避开父色. 完全图有 $q(q-1)\cdots(q-n+1)$. 对 $n\ge3$ 的圈, 删除一边得到树, 收缩得到更短的圈; $C_3$ 可直接计算, 归纳得
\[
P_{C_n}(q)=(q-1)^n+(-1)^n(q-1).
\]
这是计数恒等式, 因在所有非负整数成立, 也作为多项式恒等式在任意实数或复数成立.
\originalsection{容斥与断圈}
令 $c(A)$ 为生成子图 $(V,A)$ 的分支数, 包括孤立点. 对每条边的坏事件“端点同色”做容斥, 同时要求 $A$ 中边端点同色时, 每个分支只能选一个公共颜色, 得
\begin{equation}\label{eq:chromaticsubset}
P_G(q)=\sum_{A\subseteq E}(-1)^{|A|}q^{c(A)}.
\end{equation}
\begin{theorem}[Whitney 断圈定理]\label{thm:nbc}
给简单图的边固定全序. 一个断圈是从某个圈删去其最大边所得的边集. 令 $b_i$ 为不包含任何断圈的 $i$ 边子集数, 则 $P_G(q)=\sum_i(-1)^i b_i q^{n-i}$. 特别地, 系数按符号交替, $b_0=1$, $b_1=m$.
\end{theorem}
\begin{proof}
在容斥和中, 若 $A$ 含断圈, 存在一条边 $e$ 的两端可由 $A$ 中小于 $e$ 的边连接. 称其为可切换边, 选最大的这种 $e$, 将 $A$ 改为 $A\mathbin\triangle\{e\}$. 加删 $e$ 都不改变分支数, 而使符号相反.

须检查选择没有改变. 小于或等于 $e$ 的边的可切换性不依赖 $e$ 本身. 对大于 $e$ 的候选边 $f$, 若一条连接路使用 $e$, 可用原先小于 $e$ 的连接路替换它, 所以加入或删除 $e$ 不改变由小于 $f$ 边产生的可达关系. 因而最大可切换边仍为 $e$, 操作是无固定点的对合, 这些项两两消去.

剩余边集不含断圈, 特别不能含圈, 因而为森林, $c(A)=n-|A|$. 按边数归类得公式. 简单图中没有一边断圈, 所以每条单边集都计入 $b_1$.
\end{proof}
原手写稿把断圈定义为删去最小边; 将全序反转就与这里的版本一致. 两种约定不能在同一证明中混用.
\originalsection{无圈定向}
\begin{theorem}[Stanley 的特殊取值]\label{thm:stanley}
无自环图 $G$ 的无有向圈定向数 $a(G)$ 满足 $a(G)=(-1)^nP_G(-1)$.
\end{theorem}
\begin{proof}
对任一非自环边 $e=uv$, 考察 $G-e$ 的无圈定向. 若已有从 $u$ 到 $v$ 的有向路, 新边只能定向为 $u\to v$; 反方向会形成圈. 若已有反向路则同理. 无圈性排除两种路同时存在. 若两方向均不可达, 新边有两种选择.

这种两端不可达的定向与 $G/e$ 的无圈定向对应: 识别 $u,v$ 后不会新增有向圈, 否则识别前有一条两端连接路. 若有共同邻点, 两条边的方向必须一致指入或指出, 否则产生 $u$--$v$ 有向路, 所以重复边的方向相容. 反过来展开识别点也得到唯一的这种定向. 因而 $a(G)=a(G-e)+a(G/e)$. 有向多重图中一对端点相同、方向相反的弧组成长度2的有向圈. 因而平行边在无圈定向中必须同向, 可合并; 自环使 $a=0$. 无边图有唯一空定向. 将染色删除收缩式乘以 $(-1)^n$, 因收缩少一点, 得相同的加法递推与基值, 归纳即得.
\end{proof}
\begin{example}
求 $C_4$ 的染色多项式与无圈定向数.
\end{example}
\begin{solution}
$P_{C_4}(q)=(q-1)^4+(q-1)=q^4-4q^3+6q^2-3q$. 代入 $q=-1$ 得14. 也可直接看16种边方向中, 恰有顺时针与逆时针两种形成有向圈, 所以剩14种.
\end{solution}
\originalsection{Tutte 多项式}
顶点染色只是许多边子集计数之一. 为同时记录“还欠多少连接”和“已经多出多少圈”, 定义秩 $r(A)=n-c(A)$, 零度 $|A|-r(A)$. 二者分别是该生成子图的生成森林边数和多余边数.
\begin{definition}
允许自环和平行边的有限图的 Tutte 多项式为
\[
T_G(x,y)=\sum_{A\subseteq E}(x-1)^{r(E)-r(A)}(y-1)^{|A|-r(A)}.
\]
两个指数都是非负整数, 无边图的值为1.
\end{definition}
\begin{theorem}\label{thm:tuttedel}
普通非桥非自环边满足 $T_G=T_{G-e}+T_{G/e}$. 桥满足 $T_G=xT_{G/e}$, 自环满足 $T_G=yT_{G-e}$.
\end{theorem}
\begin{proof}
把子集分为含 $e$ 和不含 $e$ 两类. 对普通边, 删边不改变 $r(E)$; 不含边的部分直接为删除项. 含边写为 $B\cup\{e\}$, 缩边后秩变为 $r_G(B\cup\{e\})-1$, 全图秩也减1, 零度不变, 得收缩项.

若 $e$ 为桥, 令 $H=G/e$. 对 $B\subseteq E\setminus\{e\}$, 有 $r_G(B)=r_H(B)$、$r_G(B\cup\{e\})=r_H(B)+1$, 全图秩则满足 $r_G(E)=r_H(E(H))+1$. 所以不含桥的项多出 $(x-1)$, 含桥的项对应因子1, 两者相加为 $x$. 若 $e$ 为自环, 加它不改变秩, 却使零度加1, 因子为 $1+(y-1)=y$.
\end{proof}
\begin{proposition}\label{prop:tuttevalues}
设 $G$ 连通. $T_G(1,1)$ 是生成树数, $T_G(2,1)$ 是森林边子集数, $T_G(1,2)$ 是连通生成子图数, $T_G(2,2)=2^m$. 还满足
\[
P_G(q)=(-1)^{n-c(G)}q^{c(G)}T_G(1-q,0),\qquad a(G)=T_G(2,0).
\]
\end{proposition}
\begin{proof}
在 $(1,1)$ 处, 只有两个指数都为0的子集保留, 即连通且无圈的生成树; $0^0$ 在多项式代入时按常数1理解. 在 $(2,1)$ 处只保留零度0的森林; 在 $(1,2)$ 处只保留秩达到全秩的子集; 在 $(2,2)$ 处每子集贡献1.

代入 $(1-q,0)$ 后, 子集项为 $(-q)^{r(E)-r(A)}(-1)^{|A|-r(A)}$. 乘前置因子使总符号为 $(-1)^{|A|}$, $q$ 指数为 $n-r(A)=c(A)$, 与容斥式一致. 再用 Stanley 公式取 $q=-1$, 得 $a=T_G(2,0)$.
\end{proof}
\originalsection{生成树的活动性}
\begin{definition}
固定连通图边的全序及生成树 $T$. 非树边 $e$ 加到树中形成的唯一圈为基本圈; 若 $e$ 为其中最大边, 称它外部活动. 树边 $f$ 删去后产生的两部分之间所有原边构成基本割; 若 $f$ 为其中最大边, 称它内部活动. 记活动边数为 $e(T),i(T)$.
\end{definition}
\begin{theorem}\label{thm:activities}
$T_G(x,y)=\sum_{T\text{ 生成树}}x^{i(T)}y^{e(T)}$.
\end{theorem}
\begin{proof}
对边数归纳, 每次选全序中最小边 $h$. 若它是桥, 每棵生成树都含它, 基本割只有自身, 因而内部活动; 缩去它后其他活动判定不变, 得因子 $x$. 若是自环, 每棵树都不含它, 基本圈只有自身, 外部活动, 得因子 $y$.

若 $h$ 普通, 它在任何生成树中都不活动: 含它时基本割有另一边, 不含时基本圈有另一边, 而 $h$ 是最小的. 不含 $h$ 的生成树对应 $G-h$ 的生成树; 含 $h$ 的对应 $G/h$ 的生成树. 对其他边, 基本圈或基本割在删除收缩下仅可能丢去最小边 $h$, 或将经过 $h$ 的连通结构识别, 最大边身份不改变. 具体地, 树中其他边的基本割对应切开同样的两部分, 非树边的基本圈对应收缩同一路段; 删除项只是从基本割候选中去掉 $h$. 因而其他边活动性全保留. 两类求和与 Tutte 普通边递推相同. 无边单顶点为基值1, 归纳完成.
\end{proof}
这也解释系数为什么非负, 尽管子集表达式中有 $(x-1),(y-1)$. 活动性依赖所选全序, 而整个和与全序无关. 对圈 $C_n$, 递推给 $T_{C_n}=x^{n-1}+x^{n-2}+\cdots+x+y$, 在 $(1,1)$ 处得到 $n$ 棵生成树.

对连通平面图还有 $T_{G^*}(x,y)=T_G(y,x)$. 证明可直接用删缩递推: 原边的删除对应对偶边收缩, 原边收缩对应对偶边删除, 桥与自环交换. 桥与自环用交换后的乘法递推, 不对删桥后的非连通图调用这里的对偶定义. 单顶点无边图及其对偶的值均为1, 归纳即得. 因而树数在对偶下相等, 与命题\ref{prop:dual-tree}的双射一致.
\originalsection{习题与解答}
\begin{exercise}\label{ex:poly1}
求含 $n$ 点、$c$ 个分支的森林的染色多项式和 Tutte 多项式.
\end{exercise}
\begin{exercise}\label{ex:poly2}
证明 $n$ 个标号顶点上的连通简单图数 $c_n$ 满足 $c_1=1$ 及
\[
c_n=2^{\binom n2}-\sum_{j=1}^{n-1}\binom{n-1}{j-1}c_j2^{\binom{n-j}2}.
\]
\end{exercise}
\Needspace{4\baselineskip}
习题\ref{ex:poly1}的解答.
\begin{solution}
各树独立染色, 每个根有 $q$ 种选择, 其余点有 $q-1$ 种, 得 $q^c(q-1)^{n-c}$. 全部 $n-c$ 条边为桥, 所以 $T=x^{n-c}$. 孤立点不增加 Tutte 因子, 但增加染色多项式的 $q$ 因子.
\end{solution}
习题\ref{ex:poly2}的解答.
\begin{solution}
全部图数为 $2^{\binom n2}$. 按包含标号1的连通分支的大小 $j$ 分类: 另选 $j-1$ 点有 $\binom{n-1}{j-1}$ 种, 内部为一个连通图有 $c_j$ 种, 其余点任意成图有 $2^{\binom{n-j}2}$ 种, 两部分之间没有边. 对 $j=1,\ldots,n$ 求和, 把 $j=n$ 项移到左边. 这是原课完全图 Tutte 取值的一个可独立复算的计数实例.
\end{solution}
